%%% MK: Vorschlag
\section{Conclusion and Future Work}
\begin{comment}
This pattern collection presents simple solutions that make it easier for blind and visually impaired people to access graphically represented technical content with little effort. 

%The solutions were is paper has presented and empirically evaluated collection of patterns designed to make graphical representations in computer science and technical informatics education accessible to blind and visually impaired students.

The proposed solutions -- PlantUML or YAML formats for UML class diagrams, simplified LTspice netlists for electronic circuits, Java-based templates for logic circuits, and structured Markdown for \ac{kv} maps -- were derived from four years of practical experience with blind students in programming and technical informatics courses.

Structured interviews confirmed that these approaches are comprehensible, practical, and pedagogically effective, significantly lowering accessibility barriers and fostering inclusion in mixed-ability classrooms. Notably, the extracted textual content has proven valuable not only to visually impaired learners but also to sighted students, who appreciate the concise and well-structured representations. Our approach helps to create easy accessible course content using open source tools, ensuring equitable participation for blind, visually impaired, and sighted students alike.

Building on this foundation, future work will explore the integration of generative artificial intelligence (AI) to further enhance accessibility.
Our current work includes the automated generation of textual descriptions for graphical content and the creation of audio podcasts from course materials and textbooks to enable multi-modal access for a broader range of learners.

Generative AI for accessible content production has recently gained significant attention as a promising strategy for reducing educational barriers.
Research demonstrates that large language models and image-to-text systems can automatically generate alternative text, audio descriptions, or multi-modal summaries, thereby supporting blind and visually impaired learners in \ac{stem}  and other fields \cite{Karamolegkou2025}, \cite{Schauer2025}.

Upcoming research will therefore expand the existing pattern collection with new solutions leveraging AI-driven methods, investigate hybrid approaches that combine text-based formats with haptic or auditory feedback. Findings will be used to extend this pattern collection.
\end{comment}


Generative artificial intelligence has emerged as a promising approach to improving the accessibility of educational content and reducing barriers in higher education. Recent research indicates that large language models and image-to-text systems can support the automated generation of alternative text, audio descriptions, and multimodal summaries, thereby enhancing access to learning materials for \ac{bvi} learners in \ac{stem}  and related disciplines \cite{Karamolegkou2025,Schauer2025}. These developments highlight the potential of AI-assisted content creation as a scalable complement to established accessibility practices.

Building on these insights, future work will extend the existing pattern collection with solutions that incorporate AI-driven methods for accessible content production. In addition, we will explore hybrid approaches that combine text-based representations with haptic or auditory feedback to address diverse accessibility needs. The results will be systematically evaluated and used to further refine and expand the proposed pattern collection.

\subsection*{Acknowledgements} 

